\section{Exact integral algebra of the printed monodromy data}

This section replaces opaque computational status reports by certificates
that can be checked by hand.  Let
\[
 V=\mathbb Z\langle\gamma,u,w,\delta\rangle,\qquad
 \Lambda=V^*
 =\mathbb Z\langle\widehat\gamma,\widehat u,
                         \widehat w,\widehat\delta\rangle .
\]
Vectors are columns in the displayed ordered bases.  The printed monodromy on
$V$ is
\[
T_1=
\begin{pmatrix}
1&0&-6&2\\0&-1&1&1\\0&-1&0&1\\0&0&0&1
\end{pmatrix},
\qquad
T_2=
\begin{pmatrix}
1&6&0&-3\\0&0&-1&1\\0&1&0&0\\0&0&0&1
\end{pmatrix}.
\tag{3.1}
\]
Put $T_0=(T_1T_2)^{-1}$ and let
$A_j=(T_j^{-1})^t$, $M_0=(T_0^{-1})^t$ act on $\Lambda$.

\begin{proposition}[monodromy certificate]\label{prop:monodromy-certificate}
\Verified\ Direct integral multiplication gives
\[
T_1^3=I,\quad T_2^4=I,\quad
T_0=
\begin{pmatrix}
1&0&0&1\\0&1&-1&0\\0&0&1&0\\0&0&0&1
\end{pmatrix},
\quad (T_0-I)^2=0,
\]
and
\[
A_1=
\begin{pmatrix}
1&0&0&0\\6&0&1&0\\-6&-1&-1&0\\-2&1&0&1
\end{pmatrix},
\quad
A_2=
\begin{pmatrix}
1&0&0&0\\0&0&-1&0\\-6&1&0&0\\3&0&1&1
\end{pmatrix},
\quad
M_0=
\begin{pmatrix}
1&0&0&0\\0&1&0&0\\0&1&1&0\\-1&0&0&1
\end{pmatrix},
\]
with $A_1A_2M_0=I$.
\end{proposition}

\begin{proof}
Multiplying the first two matrices in (3.1) gives
\[
T_1T_2=
\begin{pmatrix}
1&0&0&-1\\0&1&1&0\\0&0&1&0\\0&0&0&1
\end{pmatrix}.
\]
Its inverse is the displayed $T_0$.  Thus
$T_0-I$ sends $(a,b,c,d)^t$ to $(d,-c,0,0)^t$, and applying it twice
gives zero.  The complete order certificate is
\[
T_1^2=
\begin{pmatrix}
1&6&-6&-2\\0&0&-1&1\\0&1&-1&0\\0&0&0&1
\end{pmatrix},
\quad T_1^2T_1=I,
\]
\[
T_2^2=
\begin{pmatrix}
1&6&-6&0\\0&-1&0&1\\0&0&-1&1\\0&0&0&1
\end{pmatrix},
\quad
T_2^3=
\begin{pmatrix}
1&0&-6&3\\0&0&1&0\\0&-1&0&1\\0&0&0&1
\end{pmatrix},
\quad T_2^3T_2=I.
\]
In particular $T_1\ne I$ and $T_2^2\ne I$.  Transposing the integral inverses gives
the three displayed contragredient matrices; finally
$A_1A_2=(M_0)^{-1}$ by the definitions, hence $A_1A_2M_0=I$.
\end{proof}

\subsection{The cusp invariant, image, and coinvariant lattices}

\begin{theorem}[all exterior-power cusp lattices]
\label{thm:cusp-lattices}
\Verified\ The equality
\[
 \ker(M_0-I)=\operatorname{im}(M_0-I)
 =\mathbb Z\langle\widehat w,\widehat\delta\rangle
\tag{3.2}
\]
holds as an equality of saturated sublattices of $\Lambda$.  On $V$,
\[
V^{T_0}=\mathbb Z\langle\gamma,u\rangle,
\]
\[
(\textstyle\bigwedge^2V)^{T_0}
=\mathbb Z\langle
\gamma u,\gamma\delta,uw,\gamma w-u\delta\rangle ,
\qquad
(\textstyle\bigwedge^3V)^{T_0}
=\mathbb Z\langle\gamma uw,\gamma u\delta\rangle .
\tag{3.3}
\]
Every displayed invariant lattice is saturated.  Consequently the invariant
ranks in degrees $0,1,2,3,4$ are $(1,2,4,2,1)$.
\end{theorem}

\begin{proof}
From the displayed matrix,
\[
(M_0-I)(a,b,c,d)^t=(0,0,b,-a)^t.
\]
Its kernel is given by $a=b=0$, and its image consists of every vector
$(0,0,b,-a)^t$; this proves (3.2).  The two generators are members of the
chosen basis, so the lattice is saturated.

The action of $T_0$ on basis vectors is
\[
T_0\gamma=\gamma,\quad T_0u=u,\quad
T_0w=w-u,\quad T_0\delta=\delta+\gamma.
\tag{3.4}
\]
Thus a vector $a\gamma+bu+cw+d\delta$ is fixed exactly when $c=d=0$.
For degree two, write
\[
x=a\gamma u+b\gamma w+c\gamma\delta
  +d\,uw+e\,u\delta+f\,w\delta.
\]
Expanding (3.4) and collecting coefficients gives
\[
(\textstyle\bigwedge^2T_0-I)x
=(-b-e+f)\gamma u-f\gamma w-fu\delta .
\]
Hence $f=0$ and $e=-b$, with $a,b,c,d$ arbitrary, which is exactly the
middle lattice in (3.3).  These are integral equations with a unit
coefficient, so their solution lattice is saturated.

In the ordered basis
$(\gamma uw,\gamma u\delta,\gamma w\delta,uw\delta)$,
\[
(\textstyle\bigwedge^3T_0-I)(a,b,c,d)^t
=(d,-c,0,0)^t.
\]
The fixed lattice is therefore generated by the first two basis vectors and
is saturated.  Degrees zero and four are fixed because $\det T_0=1$.
\end{proof}

\subsection{Global coinvariants}

Let $G=\langle T_1,T_2\rangle$ and, for a $G$-lattice $L$, write
$L_G=L/((T_1-I)L+(T_2-I)L)$.

\begin{theorem}[coinvariant certificates]\label{thm:coinvariants}
\Verified\ The common invariant lattices are
\[
 V^G=\mathbb Z\gamma,\qquad
 (\textstyle\bigwedge^2V)^G=\mathbb Zq,\qquad
 (\textstyle\bigwedge^3V)^G=\mathbb Z\gamma uw,\qquad
 (\textstyle\bigwedge^4V)^G=\mathbb Z\operatorname{vol},
\]
where $q=uw+6\gamma\delta$ and
$\operatorname{vol}=\gamma uw\delta$.  There are torsion-free
identifications
\[
V_G=\mathbb Z[\delta],\qquad
(\textstyle\bigwedge^2V)_G=\mathbb Z[\gamma\delta],
\qquad
(\textstyle\bigwedge^3V)_G=\mathbb Z[uw\delta],
\tag{3.5}
\]
and $[uw]=6[\gamma\delta]$ in degree two.
\end{theorem}

\begin{proof}
Solving $(T_1-I)v=0$ and $(T_2-I)v=0$ in (3.1) gives
\[
V^{T_1}=\langle\gamma,2u+w+3\delta\rangle,\qquad
V^{T_2}=\langle\gamma,u+w+2\delta\rangle.
\]
If a vector belongs to both lattices, comparison of its $u$- and
$w$-coordinates gives simultaneously $2r=t$ and $r=t$ for the two
non-$\gamma$ coefficients; hence $r=t=0$.  Thus $V^G=\mathbb Z\gamma$.

The higher common invariants follow from the same integral nullspace
calculation.  In the ordered bases
\[
 (\gamma u,\gamma w,\gamma\delta,uw,u\delta,w\delta),
 \qquad
 (\gamma uw,\gamma u\delta,\gamma w\delta,uw\delta),
\]
the simultaneous equations
$(\bigwedge^kT_1-I)x=(\bigwedge^kT_2-I)x=0$ are respectively equivalent to
\[
 x_1=x_2=x_5=x_6=0,\qquad x_3=6x_4,
\]
and
\[
 x_2=x_3=x_4=0.
\]
Thus the integral solution lattices are generated by
$uw+6\gamma\delta=q$ and $\gamma uw$.  Both generators are primitive.
Finally $\det T_1=\det T_2=1$, so the common degree-four invariant lattice
is $\mathbb Z\operatorname{vol}$.

In $V_G$, columns of $T_1-I$ and $T_2-I$ give, among others,
\[
u+w=0,\qquad u-3\gamma=0,\qquad 2\gamma+u+w=0,
\qquad 2u+w=0.
\]
The first two relations give $w=-u$ and $u=3\gamma$; the last two then give
$2\gamma=0$ and $3\gamma=0$.  Hence $\gamma=u=w=0$.  Every column has
zero $\delta$-coordinate, so no relation involves $\delta$.  This proves
$V_G=\mathbb Z[\delta]$, including absence of torsion.

For degree two use the ordered basis
\[
(\gamma u,\gamma w,\gamma\delta,uw,u\delta,w\delta).
\]
In this basis the twelve columns of the relation map are
\[
B_2=
\left[
\begin{array}{rrrrrr|rrrrrr}
-2&1&1&-6&2&-8&-1&-1&1&-6&6&-3\\
-1&-1&1&-6&2&-6&1&-1&0&0&3&0\\
0&0&0&0&0&-6&0&0&0&0&6&0\\
0&0&0&0&0&1&0&0&0&0&-1&0\\
0&0&0&0&-2&1&0&0&0&0&-1&-1\\
0&0&0&0&-1&-1&0&0&0&0&1&-1
\end{array}
\right],
\tag{3.6a}
\]
where the two blocks are $\bigwedge^2T_1-I$ and
$\bigwedge^2T_2-I$.  Its column lattice is
\[
\operatorname{im}
[\,\textstyle\bigwedge^2T_1-I\mid\bigwedge^2T_2-I\,]
=\langle\gamma u,\gamma w,u\delta,w\delta,
          uw-6\gamma\delta\rangle .
\tag{3.6}
\]
For completeness, the third columns of the two blocks yield
$\gamma u$ and $\gamma u+\gamma w$; their difference yields $\gamma w$.
After these are removed, the fifth and sixth columns give
$2u\delta+w\delta$ and $u\delta+w\delta$, hence both
$u\delta$ and $w\delta$; the remaining nonzero relation is
$uw-6\gamma\delta$.  Conversely, substituting the twelve displayed columns
of the two exterior powers shows that they all lie in the right side.
After ordering the ambient basis as
$(\gamma u,\gamma w,u\delta,w\delta,uw,\gamma\delta)$,
the five generators have matrix $I_4\oplus(1,-6)^t$.
Its maximal minors have gcd $1$, so (3.6) is saturated.  The quotient is
therefore freely generated by $[\gamma\delta]$, with
$[uw]=6[\gamma\delta]$.

For degree three, let
$(a,b,c,d)=(\gamma uw,\gamma u\delta,\gamma w\delta,uw\delta)$.
The full relation matrix is
\[
B_3=
\left[
\begin{array}{rrrr|rrrr}
0&0&1&2&0&-1&0&-3\\
0&-2&1&-6&0&-1&-1&-6\\
0&-1&-1&-6&0&1&-1&0\\
0&0&0&0&0&0&0&0
\end{array}
\right].
\tag{3.6b}
\]
If $c_{i,r}$ denotes column $r$ of the $T_i$ block, then
$c_{1,2}-c_{2,3}=-b$, $c_{2,3}=-b-c$, and
$c_{1,3}=a+b-c$.  Hence $a,b,c$ belong to the column lattice.
Conversely every displayed column lies in $\langle a,b,c\rangle$, and the
fourth row is zero.  The relation lattice is therefore exactly the primitive
summand $\langle a,b,c\rangle$, so the quotient is
$\mathbb Z[uw\delta]$ with no hidden torsion.
\end{proof}

\subsection{Finite quotient fibres}

Put
\[
\eta_1=2u+w+3\delta,\qquad
\eta_2=u+w+2\delta,\qquad
\mathbf q=uw+6\gamma\delta,
\]
$m_1=3$, $m_2=4$, and
$v_1=(1,2,-4,0)^t$, $v_2=-(1,3,-3,0)^t$ in $\Lambda$.

\begin{proposition}[finite-fibre pullback indices]
\label{prop:finite-indices}
\Verified\ For the free affine actions and finite fibres geometrically
realized in \cref{prop:candidate-finite-fillings}, one has
\[
S_j=A_0^{(j)}/\langle x\mapsto A_jx+v_j/m_j\rangle
\]
with $H^*(A_0^{(j)};\mathbb Z)=\bigwedge^*V$, the
degree-$m_j$ covering $\pi_j:A_0^{(j)}\to S_j$ has injective pullback with
indices, in degrees $0$ through $4$,
\[
(1,3,1,1,3),\qquad (1,4,2,2,4).
\tag{3.7}
\]
Moreover
\[
\pi_1^*H^1(S_1)=\langle3\gamma,\eta_1\rangle,\qquad
\pi_2^*H^1(S_2)=\langle4\gamma,\eta_2\rangle,
\]
\[
\pi_1^*H^2(S_1)=(\textstyle\bigwedge^2V)^{T_1},\qquad
\pi_2^*H^2(S_2)
=\{a\gamma\eta_2+b\mathbf q:a\equiv b\pmod2\}.
\tag{3.8}
\]
The identification of these abstract lattices with the actual finite-fibre
stalks is \cref{lem:finite-filling-stalk-map}; hence no geometric hypothesis
remains open in this proposition.
\end{proposition}

\begin{proof}
The fundamental group is the extension
\[
\Gamma_j=\langle\Lambda,g\mid
g\lambda g^{-1}=A_j\lambda,\ g^{m_j}=v_j\rangle .
\]
After abelianization,
\[
H_1(S_j)=
(\Lambda_{A_j}\oplus\mathbb Zg)/
\mathbb Z([v_j],-m_j).
\tag{3.9}
\]
The matrices $A_j-I$ have Smith form $\operatorname{diag}(1,1,0,0)$,
so $\Lambda_{A_j}$ is free of rank two.  The relation in (3.9) is
primitive: if a prime divides $m_j$, its image $[v_j]$ modulo that prime is
nonzero because $\gamma((A_j-I)\Lambda)=0$ and
$\gamma(v_1)=1$, $\gamma(v_2)=-1$.  Thus $H_1(S_j)$ is free of rank two.
The quotient is a closed oriented four-manifold (the linear parts have
determinant one) and
$e(S_j)=e(A_0^{(j)})/m_j=0$.  Poincar\'e duality and the universal
coefficient theorem therefore give
\[
H_k(S_j;\mathbb Z)=(\mathbb Z,\mathbb Z^2,
\mathbb Z^2,\mathbb Z^2,\mathbb Z)
\]
with no torsion.
Indeed, Poincar\'e duality gives
$H_3(S_j)\simeq H^1(S_j)\simeq\mathbb Z^2$ and
$H^3(S_j)\simeq H_1(S_j)\simeq\mathbb Z^2$.  The universal coefficient
sequence for $H^3$ then forces
$\operatorname{Ext}(H_2(S_j),\mathbb Z)=0$, so $H_2(S_j)$ is torsion-free;
Euler characteristic zero gives its rank as two.

Transfer gives $\pi_{j*}\pi_j^*=m_j$; hence $\pi_j^*$ is injective and
its image lies in the $T_j$-invariants.  The
Lyndon--Hochschild--Serre spectral sequence for
$1\to\Lambda\to\Gamma_j\to\mathbb Z/m_j\to1$
\cite[Chapter~VII, \S6]{Brown1982} has
\[
E_2^{p,q}=H^p(\mathbb Z/m_j;\textstyle\bigwedge^qV).
\]
It computes $H^*(S_j;\mathbb Z)$ because the universal cover of this affine
torus quotient is $\mathbb R^4$, hence contractible.
Its transgression on an invariant $\xi\in V$ evaluates the extension
class: $d_2\xi=\xi(v_j)\bmod m_j$.  Since
$V^{T_j}=\langle\gamma,\eta_j\rangle$ and
$\eta_j(v_j)=0$, the degree-one images and indices in (3.8) follow.

In degree two, the integral invariant lattices are
\[
(\textstyle\bigwedge^2V)^{T_j}
=\langle\gamma\eta_j,\mathbf q\rangle
\]
with no omitted finite-index overlattice.  Indeed, writing an invariant as
$(a,b,c,d,e,f)$ in the ordered basis
$(\gamma u,\gamma w,\gamma\delta,uw,u\delta,w\delta)$,
the $T_1$ equations are
$e=f=0$, $a=2b$, $c=3b+6d$, while the $T_2$ equations are
$e=f=0$, $a=b$, $c=2b+6d$.  Using $\gamma u w\delta$ as orientation,
\[
(\gamma\eta_j)^2=0,\quad
(\gamma\eta_1)\mathbf q=3,\quad
(\gamma\eta_2)\mathbf q=2,\quad
\mathbf q^2=12.
\tag{3.10}
\]
Let $L_j=\pi_j^*H^2(S_j)$.  The intersection lattice of $S_j$ is
unimodular, and pullback multiplies its form by $m_j$.  Hence
$|\det L_j|=m_j^2$.  The two invariant Gram determinants from (3.10) are
$9$ and $4$.  If $k_j$ is the index of $L_j$ in the invariant lattice,
$k_j^2|\det(\bigwedge^2V)^{T_j}|=m_j^2$, so $k_1=1$ and $k_2=2$.

For $j=2$, multiplicativity of the spectral sequence gives
$d_2(\gamma\eta_2)=-[\eta_2]\ne0$ in
$\langle\gamma,\eta_2\rangle/
\langle4\gamma,2\eta_2\rangle$.  Thus
$\gamma\eta_2\notin L_2$.  The remaining index-two possibilities are
$L_2=\langle2\gamma\eta_2,a\gamma\eta_2+\mathbf q\rangle$,
$a\in\{0,1\}$.  The covering is \'{e}tale, so
$\pi_2^*K_{S_2}$ is trivial and $K_{S_2}^{\otimes4}$ is trivial.
Because $H^2(S_2;\mathbb Z)$ is torsion-free, $c_1(S_2)=0$ and
$w_2(S_2)=0$; the intersection form is even.  Dividing the pullback
pairings by four gives
\[
\begin{pmatrix}0&1\\1&a+3\end{pmatrix},
\]
so $a+3$ is even and $a=1$.  This proves (3.8).

For clarity, the displayed differential contains no suppressed coefficient
calculation.  For the cyclic group generated by $T_2$,
\[
H^2(\mathbb Z/4;V)=V^{T_2}/(I+T_2+T_2^2+T_2^3)V
=\langle\gamma,\eta_2\rangle/
  \langle4\gamma,2\eta_2\rangle.
\]
The extension class gives $d_2\gamma=\gamma(v_2)=-1$ and
$d_2\eta_2=\eta_2(v_2)=0$; the derivation rule therefore gives
$d_2(\gamma\eta_2)=-[\eta_2]$, which is nonzero because
$\eta_2\notin\langle4\gamma,2\eta_2\rangle$.

In degree four the index is the covering degree.  For degree three, direct
exterior calculation gives the primitive invariant bases
\[
(\textstyle\bigwedge^3V)^{T_1}
=\langle\gamma uw,\,uw\delta-4\gamma u\delta-2\gamma w\delta\rangle,
\]
\[
(\textstyle\bigwedge^3V)^{T_2}
=\langle\gamma uw,\,uw\delta-3\gamma u\delta-3\gamma w\delta\rangle.
\]
Against $(\gamma,\eta_j)$ their pairing matrices are
\[
\begin{pmatrix}0&1\\-3&0\end{pmatrix},
\qquad
\begin{pmatrix}0&1\\-2&0\end{pmatrix},
\]
so the absolute determinants are $3$ and $2$.  If the degree-three
index is $r_j$, the degree-one indices already proved are $3$ and $4$.
Pullback multiplies the perfect $H^1(S_j)\times H^3(S_j)$ pairing by
$m_j$, whose rank-two determinant is $m_j^2$.  Hence
$3\cdot3r_1=3^2$ and $2\cdot4r_2=4^2$; hence $r_1=1$, $r_2=2$.
Together with degrees zero through two this proves (3.7).
\end{proof}

\subsection{The terminal presentation and degree-two map}

\begin{proposition}[terminal group, including the nonabelian step]
\label{prop:terminal-group}
\Verified\ Consider
\[
G(\ell_0,\ell_1,\ell_2)=
\langle c,x,y\mid c\text{ central},\
xy=c^{\ell_0},\ x^3=c^{\ell_1},\ y^4=c^{\ell_2}\rangle
\]
and put $p=12\ell_0-4\ell_1-3\ell_2$.  If $p\ne0$, then
$G\cong\mathbb Z/|p|\mathbb Z$; if $p=0$, then $G\cong\mathbb Z$.
For $(\ell_0,\ell_1,\ell_2)=(0,1,-1)$ one has $p=-1$ and $G=1$.
\end{proposition}

\begin{proof}
The relation $xy=c^{\ell_0}$ gives $y=x^{-1}c^{\ell_0}$.  Since $c$ is
central, $x$ and $y$ commute, so the apparently nonabelian presentation is
already abelian.  In the ordered additive generators $(x,c)$ the remaining
relations form
\[
R=\begin{pmatrix}3&-\ell_1\\4&\ell_2-4\ell_0\end{pmatrix},
\qquad \det R=-p.
\]
The greatest common divisor of all entries of $R$ is $1$ because it divides
both $3$ and $4$.  Thus its Smith form is
$\operatorname{diag}(1,|p|)$ when $p\ne0$ and
$\operatorname{diag}(1,0)$ when $p=0$.  The stated groups follow.
\end{proof}

\begin{proposition}[terminal degree-two lattice]
\label{prop:terminal-degree-two}
\Verified\ Let $A:\mathbb Z^6\to\mathbb Z^4$ have columns
\[
\begin{aligned}
c_1&=(2,-4,-2,3)^t,&c_2&=(1,-2,-2,3)^t,\\
c_3&=(3,0,-4,0)^t,&c_4&=(0,1,0,-1)^t,
\qquad c_5=c_6=0.
\end{aligned}
\]
Then $\operatorname{im}A$ is saturated of rank three,
\[
\operatorname{im}A=\ker(4,2,3,2),
\qquad
\operatorname{coker}A\cong\mathbb Z,
\tag{3.11}
\]
and $(4,2,3,2)$ is the primitive left-kernel generator.
\end{proposition}

\begin{proof}
The row $\Phi=(4,2,3,2)$ annihilates the four nonzero columns:
\[
\Phi c_1=8-8-6+6=0,\quad
\Phi c_2=4-4-6+6=0,\quad
\Phi c_3=12-12=0,\quad
\Phi c_4=2-2=0.
\]
It is primitive since the gcd of its entries is $1$, so
$\ker\Phi\subset\mathbb Z^4$ is a saturated rank-three lattice.
The $3\times3$ minor on rows $(1,2,3)$ and columns $(1,2,4)$ is $2$,
whereas the minor on rows $(1,2,4)$ and the same columns is $-3$.
Consequently the ideal generated by all maximal minors contains $(2,3)=(1)$.
This proves that $A$ has rank three and saturated image.  We have
$\operatorname{im}A\subseteq\ker\Phi$ between rank-three lattices.
The finite quotient $\ker\Phi/\operatorname{im}A$ injects into the
torsion-free group $\mathbb Z^4/\operatorname{im}A$, so it is zero.
Thus equality holds, and the primitive surjection $\Phi:\mathbb Z^4\to
\mathbb Z$ induces (3.11).
\end{proof}

\begin{proposition}[all-degree terminal homology package]
\label{prop:terminal-all-degree}
\Verified\ Let $C_\bullet$ be the degree-by-degree terminal package printed
for the glued six-manifold, and assume only the following algebraic outputs
of its displayed maps:
\begin{enumerate}[label=(\roman*)]
\item $H_0(C_\bullet)=H_6(C_\bullet)=\mathbb Z$ and the degree-one
presentation is $G(\ell_0,\ell_1,\ell_2)$ of
\cref{prop:terminal-group};
\item in degree two there is an exact sequence
$0\to\mathbb Z\to H_2^\circ\to\mathbb Z\to0$, the retained boundary
lattice kills the second quotient, and its remaining generator maps to
$p$ times a primitive generator of the first copy of $\mathbb Z$;
\item the package obeys integral Poincar\'e duality and the universal
coefficient theorem in formal dimension six, and its Euler characteristic
is $2$.
\end{enumerate}
Then, for $p\ne0$,
\[
 \bigl(H_0(C_\bullet),\ldots,H_6(C_\bullet)\bigr)
 \cong
 \bigl(\mathbb Z,\mathbb Z/|p|,\mathbb Z/|p|,
 \mathbb Z/|p|,\mathbb Z/|p|,0,\mathbb Z\bigr).
\tag{3.12}
\]
In particular, the printed terminal package has the integral homology of
$S^6$ when $p=-1$.
\end{proposition}

\begin{proof}
By \cref{prop:terminal-group}, the degree-one group is
$\mathbb Z/|p|$.  In degree two, exactness and the stated retained-boundary
map first quotient $H_2^\circ$ down to the primitive first summand and then
quotient that summand by $p$; hence $H_2(C_\bullet)=\mathbb Z/|p|$.
Both groups are finite, so their Betti numbers vanish.  Universal
coefficients and duality give
\[
 H_5\cong H^1=0,
 \qquad
 H_4\cong H^2\cong
 \operatorname{Ext}(H_1,\mathbb Z)=\mathbb Z/|p|.
\]
The Euler-characteristic identity is therefore
$2=1-b_3+1$, so $b_3=0$.  Finally
\[
 H_3\cong H^3\cong
 \operatorname{Ext}(H_2,\mathbb Z)=\mathbb Z/|p|,
\]
which proves (3.12).
\end{proof}

\paragraph{Geometric boundary.}
The calculations in this section deliberately isolate the printed lattices,
presentations, and matrices.  Their geometric realization is proved
separately: the finite-band and cusp maps are identified in
\cref{prop:actual-finite-band-maps}, the based fundamental-group calculation
is \cref{record:vankampen}, and the comparison with the intrinsic Leray
transgressions is \cref{thm:mv-leray-comparison}.
